\section{Finite-scale learning for \texorpdfstring{$\gamma>1$}{gamma > 1}}\label{sec:above}
\begin{lemma}[A non-Gaussian singular-value tail]\label{lem:spectral}
Let $B=\sigma Z$ be an $N\times N$ matrix with independent $\nu$
coordinates. If $a\geq16\sigma\sqrt{N/m}$ and $B_{>a}$ retains
exactly the singular values greater than $a$, then
\[
 \Pr\{\|B_{>a}\|_F^2>K\}\leq e^{-mK/(8\sigma^2)}.
\]
\end{lemma}
\begin{proof}
The net proof in Section~\ref{sec:target-mass} gives
\[
 \Pr\{\|B\|_{\mathrm{op}}>8\sigma\sqrt{N/m}\}
 \leq2e^{(2\log9-8)N}<1/2.
\]
The convex ball $\mathcal B=\{T:\|T\|_{\mathrm{op}}\leq a/2\}$
therefore has probability at least $1/2$. Its Frobenius distance from $B$
is
\[
 D(B)^2=\sum_j(s_j(B)-a/2)_+^2.
\]
The enlargement inequality \eqref{eq:enlargement}, after scaling by
$\sigma$, implies $\Pr\{D(B)>r\}\leq e^{-mr^2/(2\sigma^2)}$.
Since $s>a$ implies $s\leq2(s-a/2)$,
$\|B_{>a}\|_F\leq2D(B)$. Take $r=\sqrt K/2$.
\end{proof}

For later use, \eqref{eq:conditional-subgaussian} and a one-dimensional
Gaussian integration identity give
\begin{equation}\label{eq:square-moment}
 \E e^{m\xi^2/4}\leq\sqrt2.
\end{equation}
Indeed $e^{m x^2/4}=\E_G e^{\sqrt{m/2}Gx}$, with $G$ standard normal;
condition first on $G$. For $D_0$ independent coordinates, Markov's
inequality consequently gives
\begin{equation}\label{eq:sum-squares}
 \Pr\{\textstyle\sum_{j=1}^{D_0}\xi_j^2>u\}
 \leq2^{D_0/2}e^{-mu/4}\leq e^{-mu/8}
 \quad(mu\geq8D_0).
\end{equation}

\begin{lemma}[Aggregate structural tail]\label{lem:structural}
Set
\[
 a_K=[8L(\sqrt K+1)^{L-2}]^{-1},\qquad J_K=(L+3)K,
 \qquad K\geq1.
\]
If
\begin{equation}\label{eq:width-above}
 N^{\gamma-1}\geq8(d+1)/m,\qquad
 N^{\gamma-1}\geq256/(ma_K^2),
\end{equation}
then
\[
 \mu_{N,\gamma}\!\left(\bigcup_{g:\kappa_N(g)>J_K}U_{N,g}\right)
 \leq L e^{-mKN^\gamma/8}.
\]
\end{lemma}
\begin{proof}
Consider the event
\[
 \tfrac dN\|W^1\|_F^2\leq K,\quad
 N\|W^L\|_2^2\leq K,\quad
 \|(W^s)_{>a_K}\|_F^2\leq K\ (2\leq s<L).
\]
The first and last weighted norms are $N^{-\gamma}$ times sums of
$N(d+1)$ and $N$ standardized squares. Equation~\eqref{eq:sum-squares}
and the first width condition bound their tails by $e^{-mKN^\gamma/8}$.
The second width condition permits Lemma~\ref{lem:spectral} at each
internal layer. A union bound bounds failure of the event by
$Le^{-mKN^\gamma/8}$.

For a draw in this event, retain only internal singular values greater
than $a_K$. The retained matrices have operator norm at most $\sqrt K$;
each discarded matrix has norm at most $a_K$. If $e_s$ is the Euclidean
activation error, then $e_1=0$ and
$e_s\leq\sqrt K e_{s-1}+a_K\sqrt N$. The readout norm is at most
$\sqrt{K/N}$, so the uniform output error is at most
\[
 a_K\sum_{j=1}^{L-2}K^{j/2}\leq1/8.
\]
For $L=2$ there is no change. A strict-margin draw therefore retains
margin greater than $3/8$. Doubling its readout restores margin greater
than $1/2$ and gives energy at most
$K+(L-2)K+4K=J_K$. Every full-margin draw outside the asserted function
class must belong to the exceptional parameter event, proving the lemma.
\end{proof}

\begin{proof}[Proof of Theorem~\ref{thm:finite-main} for $\gamma>1$]
Write $A=A_L(d,C)$ and choose
\[
 K=\max\{1,\tfrac8m[A+\log(6L/\varepsilon)]\},\qquad
 T=\left\lceil\max\{C,[8(d+1)/m]^{1/(\gamma-1)},
                [256/(ma_K^2)]^{1/(\gamma-1)}\}\right\rceil.
\]
For $N\geq T$, divide the aggregate tail by the target lower bound
\eqref{eq:target-mass}. Global-margin normalization cancels and gives
\[
 \frac{\pi(\{g:\kappa_N(g)>J_K\})}{\pi(f)}
 \leq L e^{-N^\gamma(mK/8-A)}\leq\varepsilon/6.
\]
Necessarily $f$ belongs to the remaining class, since otherwise the ratio
would be at least one. Its log cardinality is at most $H_L(d,J_K)$.
Lemma~\ref{lem:finite-posterior}, with $a=\tau=\varepsilon/6$, gives
the claimed risk for
$n\geq(6/\varepsilon)[H_L(d,J_K)+\log(1/\delta)]$.

For $C\leq N<T$, the target mass gives
$-\log\pi(f)\leq AN^\gamma\leq AT^\gamma$.
Lemma~\ref{lem:target-posterior} gives the same risk if
$n\geq(6/\varepsilon)[AT^\gamma+\log(6/(\varepsilon\delta))]$.
At fixed $L,\gamma,m,M$, all of $K,T,T^\gamma$ have polynomial
majorants in $d,C,\log(1/\varepsilon)$. Majorize
$\max\{H_L(d,J_K),AT^\gamma\}$ by $p_{L,\gamma}$ and use
\eqref{eq:confidence}. This covers every $N\geq C$.
\end{proof}
